\documentclass[11pt]{article}
\usepackage[utf8]{inputenc}
\usepackage[T1]{fontenc}
\usepackage{lmodern}
\usepackage[margin=1in]{geometry}
\usepackage{amsmath,amssymb}
\usepackage{textcomp}
\usepackage[dvipsnames]{xcolor}
\usepackage{enumitem}
\usepackage{parskip}
\usepackage{titlesec}
\usepackage[round]{natbib}
\usepackage[colorlinks=true,linkcolor=blue,citecolor=blue,urlcolor=blue]{hyperref}

\titleformat{\section}{\large\bfseries}{\thesection}{1em}{}
\titleformat{\subsection}{\normalsize\bfseries}{\thesubsection}{1em}{}
\titleformat{\subsubsection}{\normalsize\bfseries\itshape}{}{0em}{}

\title{\textbf{Composition Over Size:\\ Allocating a Fixed Verification Budget Against Shortcut-Seeking Code Agents}}
\author{}
\date{}

\newcommand{\ph}{\textit{[placeholder --- fill in after experiments are run]}}

\begin{document}
\maketitle
\vspace{-2em}

\begin{quote}
\small\textit{Assembled from the six confirmed stage outputs of an \texttt{auto-research-discussion} run (seed topic B1.5), post research-taste audit (\textcolor{ForestGreen}{\textbf{PASS}}). Status: \textbf{specification / paper draft --- no experiments have been run.} All numbers below are pre-registered design targets, not results.}
\end{quote}

\begin{abstract}
Automated verifiers that gate agent-written code before it is merged, rewarded, or deployed cannot check exhaustively, so oversight pipelines already assume only a small, fixed budget of verification calls per episode. Existing quantitative treatments model the scarce verification resource as a single, homogeneous channel --- a monitor score plus an audit threshold --- so they quantify how budget \emph{size} trades off against safety, but cannot ask whether budget \emph{composition} across qualitatively different check types changes the answer, even though a pass-rate-optimizing agent can concentrate shortcuts exactly where a naively composed budget does not reach. We build a harness that meters three distinct check flavors --- random spot-checks, targeted invariance probes, and hidden combinatorial tests --- under one common cost rule, scores them against episodes with shortcuts implanted to carry machine-checkable ground truth, sweeps allocation across the flavor simplex to trace a miss-rate-versus-budget frontier, and tests whether mixed allocations retain more detection power than any pure policy once the agent can adapt to the verifier's allocation. A held-out-novelty control isolates whether any advantage reflects targeting structure or mere unfamiliarity. We report which composition minimizes undetected shortcuts, and how much survives an allocation-aware adversary.
\end{abstract}

\section{Introduction}

Automated verifiers now stand between an LLM coding agent's output and whatever comes next --- a merge, a reward signal, a deployment. As these agents write more of the code that unit tests, code review, and CI pipelines are asked to gate, one fact about that gating is easy to overlook: none of it is exhaustive. Running every conceivable check against every submission is not economically or computationally feasible, so real oversight pipelines already build around a small, fixed budget of verification calls per episode --- a handful of tests to run, a handful of audits to spend, and nothing more. That budget is not a temporary limitation to be engineered away; it is a permanent operating condition.

Yet a central question about that fixed budget remains unaddressed: given only a handful of calls to spend, how should they be split across qualitatively different kinds of check? A verifier could spend its budget resampling from the same visible test distribution the agent has already seen, probing a specific behavioral invariant the task is supposed to satisfy, or running a hidden combinatorial test the agent never observes. These are not interchangeable ways of spending the same currency --- each is blind to whole classes of shortcuts the others would catch. And the stakes are not symmetric: a coding agent that is implicitly or explicitly optimizing for an \emph{apparent} pass rate, rather than genuine correctness, has every incentive to concentrate its shortcuts precisely where the verification budget does not reach. If the composition of a fixed budget is spent naively, that concentration goes undetected by construction, not by bad luck.

The nearest existing work treats this scarce verification resource as a single, homogeneous channel rather than a decomposable portfolio. AI-control-style oversight protocols \citep{AIControl2024} route between auditing, deferring, and shutting down based on one monitoring signal over an entire deployment trajectory; a companion line \citep{ControlTax2024} converts one monitor's ROC curve and an audit budget directly into a probability of compromise. Both can say a great deal about how \emph{much} auditing a given safety level requires, but neither can express the question of \emph{which composition} of structurally different checks buys the most detection per call, because both collapse the verification resource to a single scalar signal before the composition question can even be posed. Successors in the same line widen the protocol's \emph{action} space --- resampling suspicious actions \citep{CtrlZ2025}, formalizing the audit game itself \citep{GamesForAIControl2024} --- and, from the attacker's side, already demonstrate that a monitor-aware adversary can subvert these single-channel protocols by adapting to the one signal they rely on \citep{AdaptiveAttacksMonitors2025,AttackSelectionConcentrated2026}; what none of this work varies is the composition of the checks themselves. Classical inspection-game theory, the intellectual ancestor of this framing \citep{AvenhausInspectionGames2002}, solves budget-constrained inspector-inspectee games rigorously but with one homogeneous inspection technology and no empirical instantiation against an autonomous, pass-rate-optimizing agent. A separate line of work on reward hacking \citep{InfTimeRewardHacking2024} documents the same underlying failure mode --- that stronger optimization against an imperfect, fixed proxy produces a characteristic collapse in true performance --- but mitigates it entirely from the optimizer's side, holding the verifier itself fixed rather than asking how the verifier's own scarce resource should be spent.

Turning the composition question into something measurable surfaces its own obstacles. The three natural check flavors have structurally incomparable detection profiles: a random sample only catches shortcuts on inputs it happens to draw, an invariance probe only catches shortcuts that break a specific behavioral property, and a hidden combinatorial test only catches shortcuts that fail on unusual input combinations --- so no flavor dominates across all shortcut classes, and a fair comparison first requires making structurally different check operations commensurable at all before any ranking is possible. Even once they are comparable, a fixed budget still forces an inescapable trade-off between broad coverage and targeted depth, whose shape cannot be assumed in advance and has to be measured as a swept curve rather than argued for. Worse, any allocation recommendation is only as good as its robustness to the very adversary it is meant to catch: a deployed verifier is reused across many episodes, not consulted once, so a patient attacker who observes enough outcomes can infer which calls are predictable and redistribute shortcuts accordingly --- an allocation that looks best against a naive attacker can lose that advantage entirely once the attacker adapts. And measuring any of this at all requires a ground-truth label for whether a passing solution was genuine or a shortcut, which the very budget under study cannot afford to produce through exhaustive checking of its own.

Our approach turns each of these obstacles into a specific design choice. Every check flavor is metered under one common cost rule and evaluated against the same pool of episodes, so ``qualitatively different'' checks produce numbers on one shared axis instead of an assumed common currency. Allocation is treated as a point on the full three-flavor simplex and swept across budgets, turning the coverage-versus-depth trade-off into a measured frontier rather than an asserted trade-off. Every allocation policy is tested twice --- once against an attacker with no information about the verifier's allocation, and again against an attacker that can observe and adapt to it --- so that a policy's apparent strength cannot be an artifact of facing only a naive opponent. And ground truth is obtained by constructing shortcuts with a machine-checkable witness at generation time, paired with a control arm that strips away each targeted flavor's structure while preserving its novelty, so that any measured advantage can be attributed to what a check actually does rather than merely to it being unseen.

We evaluate this design on coding tasks spanning single-issue and multi-argument programs, comparing allocations against both a naive and an adaptive attacker, alongside a reproduced single-channel auditing protocol and floor/ceiling sanity controls.

In summary, we contribute:
\begin{itemize}[leftmargin=1.4em]
\item A blind-versus-aware stress test that measures whether mixed or randomized allocations retain more detection power than any pure or single-channel policy once the agent can infer and adapt to the verifier's own allocation --- rather than reporting a ranking that only holds against a naive attacker.
\item A cost-commensurable scoring scheme that meters three structurally distinct check flavors identically, making detection profiles that are individually incomparable comparable on one shared miss-rate axis.
\item A simplex-wide frontier sweep that turns the inherent coverage-versus-depth trade-off of a fixed verification budget into a measured Pareto curve rather than an assumed or asserted ordering.
\item A construction-based ground-truth oracle paired with a targeting-structure-ablated novelty control, isolating whether a check flavor's advantage comes from what it specifically targets or merely from being held out of the agent's visible feedback.
\end{itemize}

\section{Problem Definition}

\noindent\textbf{Formal setup.} Consider a coding task $\tau$ with a known reference solution $r_\tau$, to which an agent submits a program $s$. Each episode is a tuple $e = (\tau, s, y_e, w_e)$, where $y_e \in \{\text{genuine}, \text{shortcut}\}$ is a ground-truth label determined independently of the verifier, and $w_e$ is a machine-checkable witness (a trigger predicate or a violated invariance) that substantiates $y_e$ when it is a shortcut. A verifier has access to three structurally distinct check flavors, $F = \{R, P, H\}$ --- a uniform random spot-check drawn from a held-out input pool, a targeted invariance probe built from a semantics-preserving transformation, and a hidden combinatorial test drawn from a covering array over the task's input parameters --- each metered under one common rule: one verification call equals one sandboxed execution plus one output comparison. An allocation policy $\pi$ is a point on the simplex $\Delta^2$ over $F$, rounded to integer calls at a fixed per-episode budget $B$. The object of study is $MR(\pi, B)$, the fraction of shortcut episodes for which zero of the $\leq B$ spent calls flag a failure, evaluated both when the agent has no information about $\pi$ (allocation-blind) and when it can observe or infer $\pi$ from past outcomes (allocation-aware). The goal is to identify, for each swept $B$, which $\pi$ minimizes $MR(\pi, B)$ --- reported both pooled and broken out per shortcut-generating attack family, under a call-count cost meter as the primary axis, with wall-clock cost as a pre-committed fallback axis if the two meters disagree beyond a stated fairness bound.

\noindent\textbf{Estimand.} The quantity this paper measures is $MR(\pi, B)$ --- the fraction of ground-truth-shortcut episodes with zero flagged checks --- as a function of allocation policy $\pi$ and budget $B$, over the population of (task, attack-family) episodes drawn from the task suites described in Experiments. The unit of analysis is one episode; the decision this quantity feeds is which allocation policy (or mixture) a budget-constrained verifier designer should commit to at a given call budget and adversary-awareness regime. This estimand is stated once here and is not re-derived per research question --- every formal object introduced above and every research question in Experiments instantiates it.

\noindent\textbf{Why this is hard.} The three flavors in $F$ have non-nested detection profiles: a spot-check only catches shortcuts on the inputs it happens to draw, a probe only catches shortcuts that violate the specific invariant it checks, and a combinatorial test only catches shortcuts that fail on the particular input combinations it covers --- so no flavor dominates across all shortcut classes, and comparing them at all first requires a scheme that makes structurally different check operations commensurable on one axis. Even once commensurable, a fixed $B$ forces an inescapable trade-off between spending calls broadly versus spending them on a targeted flavor, whose shape is not derivable in advance and must be measured as a curve over swept budgets and allocations rather than assumed. The object of study is further complicated because $\pi$ is not evaluated in a vacuum: an allocation that appears best against an agent with no information about $\pi$ may not remain best once the agent can condition on past check outcomes and adapt its shortcut placement --- so any claim about the ``best'' $\pi$ is only meaningful relative to a stated assumption about what the agent can infer. Finally, $y_e$ itself cannot be produced by the budgeted verifier under study without contradiction --- computing a ground-truth label would require exactly the exhaustive checking the fixed budget is designed to avoid --- so $y_e$ must be established by a mechanism entirely outside the game being measured, and that substitution's fidelity to shortcuts an agent might discover unprompted (rather than ones deliberately implanted) is itself an assumption that has to be stated and probed, not assumed away.

\section{Method Overview}

At a conceptual level, the study is a single measurement pipeline with four cooperating parts, all operating over the same stream of coding-task episodes and the same three check flavors. A construction step first produces episodes whose genuine-vs-shortcut label is known without needing the budgeted verifier to determine it, and pairs each targeted flavor with a structure-ablated control so that any later advantage can be traced to what a flavor actually does rather than merely to novelty. A scoring step meters all three flavors under one common cost rule and evaluates them against the same episode pool, turning individually incomparable detection profiles into numbers on one shared axis. A sweep step then treats allocation as a point on the full flavor simplex and traces how the resulting miss rate moves as both the budget and the allocation vary, turning the coverage-versus-depth trade-off into a measured frontier. Finally, a stress-test step re-runs every allocation policy from that frontier against an attacker that can observe and adapt to the verifier's own allocation, checking whether a policy's apparent strength survives contact with exactly the kind of adversary the whole exercise is meant to defend against.

The \textbf{construction oracle with novelty control} implants shortcuts with a machine-checkable witness at generation time and additionally produces a targeting-structure-ablated variant of each targeted flavor, so ground truth and the structure-versus-novelty attribution are both settled before any policy is scored. The \textbf{cost-commensurable flavor scorer} takes every episode and every flavor and produces one miss-rate number per (flavor, budget, attack-family) cell under the shared cost meter. The \textbf{simplex frontier sweep} takes those per-flavor numbers and extends them across the full allocation simplex and a range of budgets, producing the miss-rate-versus-budget frontier and identifying candidate winning allocations. The \textbf{blind-versus-aware stress test} takes the frontier's candidate allocations --- including a reproduced single-channel auditing baseline --- and re-evaluates each one against an attacker with and without information about the allocation, reporting how much of each policy's advantage survives that adaptive pressure.

\section{Methodology}\label{sec:methodology}

\subsection{Construction Oracle with Novelty Control}
\textbf{Challenge.} Computing a miss rate requires a ground-truth genuine-vs-shortcut label for every episode, but the budgeted verifier being studied cannot afford the exhaustive check that would produce that label without contradicting the premise of the study.

\textbf{Insight.} Ground truth can be constructed outside the game entirely --- but that construction only earns validity if it isn't silently doing the explanatory work that should be attributed to a check flavor's actual targeting logic, which is exactly what a novelty-ablated control is built to rule out.

\textbf{Method.} Shortcuts are implanted across three attack families --- hard-coded special-casing of visible test inputs, input-conditional backdoors that fail on a rare trigger predicate, and invariance-breaking implementations correct only on canonical input formats --- each producing a machine-checkable witness at generation time, so every episode carries a label without requiring any check to determine it. In parallel, each targeted flavor (the invariance probe and the hidden combinatorial test) is paired with an ablated variant that draws checks from the identical held-out input universe but strips out the targeting logic itself --- random input mutations in place of semantics-preserving transforms, random parameter tuples in place of covering-array structure --- judged against the reference solution identically. Comparing a targeted flavor's miss rate against its own ablated variant isolates whether an observed advantage comes from what the flavor specifically targets or merely from drawing on inputs the agent has not seen.

\subsection{Cost-Commensurable Flavor Scorer}
\textbf{Challenge.} The three check flavors have non-nested detection profiles --- no single flavor catches every class of shortcut --- so there is no obvious common currency in which to compare them.

\textbf{Insight.} Commensurability does not have to be assumed; it can be constructed, by holding every flavor to one identical cost rule and evaluating all of them against the same episode pool.

\textbf{Method.} Every verification call, regardless of flavor, is metered as one sandboxed execution plus one output comparison. Random spot-checks draw without replacement from a held-out I/O pool; invariance probes apply one semantics-preserving transformation per call and check output consistency against the untransformed run; hidden combinatorial tests draw from a pairwise covering array over the task's input parameters, judged against the reference solution and never exposed in the agent's visible feedback. Running all three flavors under this identical meter against the same implanted-shortcut episodes, broken out by attack family, produces one miss-rate number per (flavor, budget, attack-family) cell that can be directly compared across flavors.

\subsection{Simplex Frontier Sweep}
\textbf{Challenge.} A fixed, small budget forces an inescapable trade-off between spending calls broadly and spending them on a targeted flavor, and the shape of that trade-off is not knowable in advance.

\textbf{Insight.} Rather than assuming any single allocation is correct, the trade-off can be turned into a measured Pareto frontier by treating allocation itself as a variable to be swept, not a design choice to be argued for.

\textbf{Method.} An allocation policy is parameterized as a point on the three-flavor simplex, integer-rounded at a given budget. The simplex is swept at a fixed grid resolution, alongside the three pure vertices and an adaptive policy that re-weights flavors per episode from a running per-flavor detection-yield estimate, across a range of swept budgets. For each swept point, the miss-rate-versus-budget curve and its trapezoidal area are recorded, yielding a frontier over which pure and mixed allocations can be directly compared for Pareto dominance.

\subsection{Blind-versus-Aware Stress Test}
\textbf{Challenge.} An allocation's apparent strength against a naive attacker says nothing about whether it survives a deployed setting, where a persistent verifier is reused across many episodes and a patient attacker can learn to predict it.

\textbf{Insight.} Mixed-strategy commitment is the standard defense against a best-responding opponent in inspection and security games; if that logic applies here, it should show up as smaller degradation specifically for mixed or randomized policies once the agent adapts, not as uniform degradation across all policies.

\textbf{Method.} Every candidate allocation from the frontier sweep --- plus a reproduced single-channel auditing baseline built by threshold-tuning one suspicion signal at matched call budget --- is evaluated under two attacker conditions: allocation-blind, where the attack policy has no information about the allocation, and allocation-aware, where the attack policy conditions on the verifier's past check outcomes and a descriptor of the allocation. The drop in each policy's detection performance between the two conditions is compared via a policy-by-condition interaction test, isolating which policies lose the least ground once the agent stops behaving naively.

\section{Experiments}\label{sec:experiments}

\subsection{Experimental Setup}

\textbf{Datasets.} APPS (competition-level programming problems) --- the implanted-backdoor substrate the closest prior budgeted-oversight work already uses; the pilot-stage domain, since its single-issue structure keeps the harness simple to validate first. BigCodeBench (function-level tasks with rich, multi-argument call signatures) --- the hidden-combinatorial-test flavor needs enough independent input parameters for a pairwise covering array to be non-trivial, which APPS's single-issue style under-stresses; this is the scale-up domain that actually tests the combinatorial flavor's advantage rather than assuming it. EvalPlus-extended HumanEval+/MBPP+ --- the source of the large held-out I/O pools the random-spot-check flavor and the unstructured-novelty control both need.

\textbf{Attack generation and elicitation risk.} Attack episodes are generated by prompted attack scaffolds over at least two frontier-model families plus one open-weights model. Two pre-registered mitigations address the elicitation risk that frontier models refuse backdoor-implantation prompts or produce only weak, easily caught shortcuts --- a failure that would flatten every frontier and make the flavors indistinguishable for reasons that have nothing to do with composition: a scripted, template-based implantation floor that cannot refuse (the same scripted attacks used by the instrument-validation positive control below), and an open-weights attacker fallback. Resource plan (a pre-registered planning estimate, not a measured cost): the pilot is roughly 100 tasks $\times$ 3 attack families $\times$ 5 policies $\times$ 4 budgets $\times$ 3 seeds $\approx$ 18{,}000 episodes at $\leq 8$ sandboxed calls each --- on the order of $10^5$ cheap sandboxed executions plus $10^3$--$10^4$ attack-generation LLM calls reused across policies and budgets --- with the scale-up stage sized at 4--6$\times$ the pilot.

\textbf{Baselines} (sanity-check comparators, not competitors to beat --- this is a measurement study, and no research question here is framed as ``our procedure beats baseline X''). The three pure allocation vertices are the necessary reference points every mixture is compared against. A single-channel suspicion-score-triggered auditing protocol, reproduced at matched call budget from the closest prior budgeted-oversight work \citep{AIControl2024,ControlTax2024}, is the external comparator for the composition question. Its cost is mapped onto the shared meter by a pre-registered rule: one monitor query is metered as one verification call on the primary call-count axis; because a monitor query and a sandboxed execution differ in wall-clock and dollar terms, this baseline is exactly where the two cost meters are most likely to disagree, and any disagreement beyond the stated fairness bound triggers the pre-committed wall-clock fallback axis (Problem Definition), reported alongside the call-count result. A no-verification floor ($B=0$, miss rate identically 1 by construction) and a visible-test-replay floor (all calls spent re-executing draws from the visible test set, expected miss rate near 1 since attacks are constructed to pass visible tests) anchor the frontier and validate that the attack-implantation filter itself is working. A label-permutation placebo --- ground-truth labels permuted within task, every policy re-scored through the identical evaluator --- is a null-by-construction arm: any apparent policy separation on permuted labels must vanish, and a non-null result here means the scoring pipeline measures an artifact, so it is checked before any research question is scored. An unstructured-novelty control, described above, isolates targeting structure from mere held-out novelty. Genuine episodes --- the false-alarm guard's denominator --- are reference solutions plus unprompted model solutions that pass the visible tests; a pre-registered audit of a random 100-episode genuine sample per model family checks for spontaneous test-special-casing, and if impurity reaches 2\%, false-alarm rates are reported with an impurity-adjusted bound and every flagged-genuine episode is manually adjudicated before the disqualification rule is applied.

\textbf{Metrics.} The miss rate $MR(\pi,B)$ and its frontier-AUC summary; the false-alarm rate on genuine episodes (a guard metric --- any flavor whose false-alarm rate exceeds a pre-specified bound is disqualified from being declared a winner regardless of its miss rate); the realized cost per episode against the declared budget; and the gap between an agent's apparent (visible-test) pass rate and its ground-truth success rate.

\textbf{Instrument validation.} The two targeted flavors are generated instruments --- a metamorphic mutation set ($P$) and a covering-array detector ($H$) --- and are validated before any research-question arm runs, then content-hash frozen before any attack generation. Soundness: the reference solution must pass 100\% of retained probes and checks on every retained task; probes or generated checks that flag the reference solution are dropped --- never patched per-task --- with the drop rate reported, which also screens out-of-domain covering-array inputs that would otherwise inflate false alarms. Positive control: at exhaustive check budget, each targeted flavor must detect at least 90\% of scripted attacks from its own detection class; a flavor that cannot detect its own attack class by construction is a broken instrument, and no null result involving it is interpretable --- such a failure is reported as an instrument failure, not as evidence about composition. The covering-array strength is fixed at $t=2$ (pairwise) on two grounds: empirical fault studies in the combinatorial-testing literature find that most interaction faults involve two or fewer parameters, and higher-strength arrays grow combinatorially, so a budget of $B$ calls would cover a vanishing fraction of a $t=3$ array at the budgets swept; an exploratory sensitivity check that re-generates the hidden-test pool at $t=3$ on the subset of tasks with enough independent input parameters is pre-registered alongside the primary analysis.

\textbf{Attack-family-mix sensitivity.} The three attack families are near-dual to the three flavors' detection classes, so the pooled family mix --- an experimenter choice --- mechanically shapes the miss-rate ordering. A pre-registered sensitivity analysis removes the winner-as-artifact-of-mix confound: RQ1's winner and RQ2's mixture-versus-vertex contrasts are recomputed under reweighted family mixes with each family's share set in turn to 10\%, 33\%, and 80\% (renormalizing the others), as a reanalysis of logged per-episode outcomes with no new runs. The allocation recommendation is stated unconditionally only if the winner is stable across all reweightings; otherwise it is stated mix-conditionally, anchored to the natural-episode audit for which mixes are deployment-realistic.

\textbf{Budget-stability of the recommendation.} The stated decision --- commit to an allocation at a given call budget --- is only as good as the winner's stability along the frontier, so a pre-registered budget-stability check recomputes RQ1's winning flavor and RQ2's mixture-versus-vertex contrasts separately at each swept budget, as a reanalysis of logged per-episode outcomes with no new runs. The allocation recommendation is stated budget-unconditionally only if the same policy wins at every swept budget; if the winner flips along the frontier, the recommendation is stated per-budget.

\textbf{Thin-cell rule.} Per-(policy, budget, condition) episode counts are logged before any research question is scored, on a cluster-adjusted effective-$n$ basis (episodes cluster by task, so raw episode counts overstate power; the design effect $\mathrm{DEFF} = 1 + (\bar m - 1)\cdot\mathrm{ICC}$ is computed with the intra-task correlation estimated from pilot data --- no value asserted in advance). Any cell that falls below the powered threshold targeted for that stage --- whether from task drops, the held-out attack-family split, or seed deduplication reducing effective episodes --- is pre-declared \textbf{inconclusive} for that cell: never treated as a null result, and never silently dropped from the reported grid. Each RQ's primary contrast is scored only over adequately powered cells; if a primary contrast's own cell is thin, that specific contrast is reported inconclusive rather than defaulting to its falsification-condition wording, which is reserved for a powered test that genuinely fails.

\textbf{Research questions:}
\begin{enumerate}[leftmargin=1.6em]
\item RQ1 --- Do targeted check flavors achieve a lower miss rate than uniform random spot-checking at matched budget against an allocation-blind attacker?
\item RQ2 --- Does any mixed allocation Pareto-dominate every pure flavor on the miss-rate-versus-budget frontier?
\item RQ3 --- Do mixed or randomized allocations retain more detection power than pure policies once the agent becomes allocation-aware?
\item RQ4 --- How does detection-power retention under an allocation-aware agent compare between a flavor-portfolio policy and a single-channel auditing baseline, at matched cost?
\item RQ5 --- Does a targeted flavor's miss-rate advantage survive when its targeting structure is ablated but its held-out novelty is preserved?
\end{enumerate}

\subsubsection{RQ1: Do targeted check flavors achieve a lower miss rate than uniform random spot-checking at matched budget against an allocation-blind attacker?}
\textbf{Setup.} Pilot stage on APPS: 2 backbones $\times$ 3 pure allocation policies $\times$ budgets $B \in \{1,2,4,8\}$ $\times$ 50 pairs of (agent, verifier) runs $\times$ 3 attack families, allocation-blind condition only. Scale-up on APPS and BigCodeBench jointly: $\geq$4 backbones $\times$ budgets $B \in \{1,2,4,8,16\}$ $\times$ $\geq$100 pairs per cell. Pairwise two-proportion z-tests between flavors within each budget, Holm-corrected across the three pairwise contrasts per budget, with a task-level cluster bootstrap (10,000 resamples) as the primary interval estimate.

\textbf{Expected outcome.} At matched budget $B \in \{4,8\}$, both the invariance probe and the hidden combinatorial test achieve a miss rate at least 10 percentage points lower than uniform random spot-checking, each contrast significant at $\alpha=0.05$ after Holm correction, with the bootstrap 95\% CI on the difference excluding zero.

\textbf{Falsification condition.} If no targeted flavor beats random spot-checking by at least 10 percentage points at any swept budget (or the bootstrap CI on the difference includes zero), the claim reverses: uniform random spot-checking is already cost-competitive with structurally targeted probing. The stronger clause --- that check-flavor composition is a second-order lever at budgets $\leq 8$ --- is licensed only if RQ2 also fails (outcome-combination rule CR-4 below), since a passing RQ2 would demonstrate that composition matters through mixing. Wrong-direction rule: if random spot-checking is instead \emph{lower} than both targeted flavors by at least 10 percentage points (same test and correction), that reversal is reported at full prominence as RQ1's headline --- targeted structure actively hurts at matched budget, and every probe or hidden-test call spent is a held-out random draw forgone.

\textbf{Hypothesis (plain language).} Implanted shortcuts are built to survive the visible test distribution, which random spot-checking draws from a similar distribution and therefore re-covers inefficiently; targeted flavors draw from regions the attack was never fit against.

\textbf{Metrics.} Per-flavor, per-budget miss rate and its bootstrap CI; false-alarm rate as a guard; per-attack-family breakdown alongside the pooled comparison; per-attacker-backbone breakdown as a pre-registered heterogeneity check --- the pooled flavor ordering is claimed as model-general only if its direction replicates within every attacker backbone, and is otherwise reported as pooled with the per-backbone divergence disclosed.

\textbf{Result.} \ph

\subsubsection{RQ2: Does any mixed allocation Pareto-dominate every pure flavor on the miss-rate-versus-budget frontier?}
\textbf{Setup.} A 10\%-grid sweep of the three-flavor simplex plus an adaptive Thompson-sampling-style re-weighting policy, evaluated at every swept budget alongside the three pure vertices; the confirmatory contrasts are two pre-specified mixture candidates --- the uniform interior mixture $(1/3,1/3,1/3)$ and the adaptive policy --- each against every pure vertex, Holm-corrected over the two candidates, with the full simplex sweep reported as exploratory (no grid mixture is selected post hoc from the same data and scored confirmatorily --- a winner's-curse guard).

\textbf{Expected outcome.} Some evaluated mixed policy achieves a frontier AUC at least 0.05 lower (roughly 5 percentage points of mean miss rate across the swept budgets) than every pure vertex's AUC, with the task-level bootstrap 95\% CI on each mixture-versus-vertex difference excluding zero.

\textbf{Falsification condition.} If no evaluated mixed or adaptive policy's frontier AUC is at least 0.05 lower than every pure vertex's (or every mixture-versus-vertex CI includes zero), the claim becomes: mixing dilutes rather than hedges --- a budget-constrained verifier should identify the single dominant flavor for its shortcut mix and commit to it, rather than split the budget.

\textbf{Hypothesis (plain language).} Because the three attack families are differentially detectable by different flavors, hedging across flavors should out-perform committing to any single one --- unless mixing merely dilutes the strongest flavor rather than covering what it misses.

\textbf{Metrics.} Frontier AUC per policy with bootstrap CI; pairwise Pareto-dominance check at each swept budget.

\textbf{Result.} \ph

\subsubsection{RQ3: Do mixed or randomized allocations retain more detection power than pure policies once the agent becomes allocation-aware?}
\textbf{Setup.} Every policy from the RQ2 frontier is re-evaluated under an allocation-aware attack condition, where the attack policy conditions on the verifier's past check outcomes and a descriptor of the allocation, using the same task set as the paired allocation-blind runs. A policy-by-condition interaction is estimated via a logistic model with task-clustered robust standard errors. A manipulation check gates the whole condition: the aware condition counts as operative only if a scripted oracle-aware attacker --- told the allocation verbatim --- degrades the best pure policy's miss rate at $B=4$ by at least 10 percentage points over blind; if this floor is not met, the awareness channel is declared inoperative and RQ3/RQ4 are reported inconclusive, an instrument failure rather than a mechanism claim.

\textbf{Expected outcome.} The best pure policy's miss-rate degradation from allocation-blind to allocation-aware at $B=4$ exceeds the best mixed or randomized policy's degradation by at least 10 percentage points, with the interaction term significant at $\alpha=0.05$.

\textbf{Falsification condition.} If the best pure and best mixed policy's blind-to-aware degradation are statistically indistinguishable (the interaction term is not significant, or the gap is under 10 percentage points), the claim becomes: allocation-awareness taxes all defenses roughly equally, refuting the mixed-strategy-robustness hypothesis for check portfolios --- the binding variable is information leakage to the agent, not budget composition.

\textbf{Hypothesis (plain language).} Mixed-strategy commitment is the standard defense against a best-responding opponent; a deterministic or near-deterministic allocation should lose more ground once it becomes predictable than a genuinely randomized one. This direction is the \emph{expected} result in the security-games literature (Related Work), so the pre-registered contribution here is its measured magnitude for executable check portfolios at matched call cost --- a refutation would be the more surprising outcome.

\textbf{Metrics.} Blind-to-aware miss-rate change per policy at each budget; interaction-term estimate and confidence interval.

\textbf{Result.} \ph

\subsubsection{RQ4: How does detection-power retention under an allocation-aware agent compare between a flavor-portfolio policy and a single-channel auditing baseline, at matched cost?}
\textbf{Setup.} The single-channel suspicion-score-triggered auditing baseline, threshold-tuned and reproduced at matched call budget, is run through the identical blind-versus-aware protocol used for the portfolio policies in RQ3, using the same interaction test, Holm-corrected jointly with RQ3's contrast since both reuse the same interaction design.

\textbf{Expected outcome.} The single-channel baseline's blind-to-aware degradation at $B=4$ exceeds the best flavor-portfolio policy's degradation by at least 10 percentage points.

\textbf{Falsification condition.} If the single-channel baseline's blind-to-aware degradation does not exceed the best portfolio policy's by at least 10 percentage points (or the interaction term is not significant), the claim becomes: a flavor-portfolio confers no additional robustness over a well-tuned single-channel auditing threshold once an agent adapts --- undercutting the case for portfolio allocation specifically as a defense against allocation-aware adversaries, independent of whatever RQ1/RQ2 find in the allocation-blind regime.

\textbf{Hypothesis (plain language).} A verification approach built around one predictable scalar signal should have less to fall back on once an agent learns to predict it, compared to a portfolio that can distribute its unpredictability across structurally different check types. That a single predictable monitor channel can be subverted by an adversary aware of it is \emph{already demonstrated empirically} for LM-monitor-based control protocols \citep{AdaptiveAttacksMonitors2025}; the pre-registered contribution here is therefore not the existence of single-channel degradation but its matched-cost magnitude relative to a check portfolio --- a comparison no prior work in that line performs.

\textbf{Metrics.} Blind-to-aware degradation for the single-channel baseline versus the best portfolio policy, same units and test as RQ3.

\textbf{Result.} \ph

\subsubsection{RQ5: Does a targeted flavor's miss-rate advantage survive when its targeting structure is ablated but its held-out novelty is preserved?}
\textbf{Setup.} Within the allocation-blind condition at $B=4$, each targeted flavor's miss rate (from RQ1) is compared against its own novelty-matched, structure-ablated variant, using a paired bootstrap over the same task set. The comparison applies only to targeted flavors that individually beat random spot-checking in RQ1, Holm-corrected over the tested flavors, with a 5-percentage-point effect floor.

\textbf{Expected outcome.} The targeted flavors' miss-rate advantage over random spot-checking shrinks to under 5 percentage points and becomes statistically indistinguishable from zero (bootstrap CI including zero) once targeting structure is ablated while held-out novelty is preserved.

\textbf{Falsification condition.} If a targeted flavor still beats its structure-ablated, novelty-matched control by 5 percentage points or more (bootstrap CI excluding zero) at $B=4$, the claim reverses: the advantage survives even with targeting logic stripped out, so it was never really about what the flavor specifically checks --- only about it being unseen --- which would undercut the case for building targeted checks at all rather than simply holding out more random inputs.

\textbf{Hypothesis (plain language).} If a targeted flavor's advantage survives even after its specific targeting logic is stripped out, the advantage was never really about \emph{what} the flavor checks --- only about it being unseen, which would undercut the case for building targeted checks at all rather than just holding out more random inputs.

\textbf{Metrics.} Miss-rate difference between each targeted flavor and its structure-ablated variant, with bootstrap CI.

\textbf{Result.} \ph

\vspace{0.5em}
\noindent\textbf{Outcome-combination rules.} The five falsification conditions above are not scored independently at verdict time: a pre-registered set of combination rules (fixed in the design specification as CR-1 through CR-7) states which claim wording is licensed under every joint outcome of RQ1--RQ5. In particular: the second-order-lever clause is licensed only when RQ1 and RQ2 both fail (CR-4); if RQ3's manipulation check fails, RQ3 and RQ4 are reported inconclusive regardless of their nominal outcomes and no awareness mechanism claim appears anywhere in the paper, with the global negative-result pivot --- triggered only if RQ1, RQ2, and an operatively-tested RQ3 all fail --- then using a pre-committed variant that drops its awareness clause (CR-6); if RQ3 passes but RQ4 fails, the portfolio's robustness value is reported as confined to comparisons among portfolio policies, not over the single-channel baseline (CR-5); and if RQ5's novelty control wins, its re-attribution (``novelty, not structure'') modifies every surviving RQ1/RQ2 claim, while RQ3's randomization claim survives restated with the re-attributed mechanism, since unpredictability is preserved under that re-attribution (CR-7).

\vspace{0.5em}
\noindent\textbf{Further evaluation.} With more time, a natural extension is to move beyond implanted shortcuts entirely and audit naturally emergent shortcut-taking behavior --- the kind an RL-trained or instruction-following agent discovers on its own rather than one deliberately implanted for the study. This is explicitly out of scope here because no affordable ground-truth oracle exists for that setting at the scale this design requires; a small validation audit against naturally occurring test-special-casing episodes is the most that is feasible without a fundamentally different evidence source.

\section{Related Work}

\textbf{Single-channel budgeted oversight.} AI-control-style protocols \citep{AIControl2024} evaluate multi-model oversight under the explicit assumption that the most capable model may be adversarial, routing between auditing, deferring, and shutting down based on one monitoring signal across an entire deployment trajectory; a companion line \citep{ControlTax2024} embeds a monitor's empirically measured ROC curve directly into an attacker-defender model that converts monitor accuracy and audit budget into a compromise probability and a safety-cost trade-off. A related sabotage-evaluation framework \citep{CtrlAltDeceit2024} extends this to explicit backdoor and generalization-failure objectives, scoring agent transcripts with LM-based suspicion monitors under both oversight-aware and oversight-unaware conditions. Successor protocols widen the \emph{action} space rather than the \emph{check} space --- resampling suspicious actions in place of a single audit-or-defer decision \citep{CtrlZ2025} --- while still routing every decision through one monitor signal. A separate, non-adversarial line combines many weak scalar verifiers into one stronger selection signal \citep{Weaver2025}, but with no adversary, no per-call budget, and no structural heterogeneity in what each verifier executes --- its estimand is selection accuracy, not adversarial miss rate. All of these treat the scarce verification resource as a single, homogeneous channel --- one monitoring signal and one threshold --- and so can characterize how much auditing a given safety level requires, but none can express the question of which composition of structurally distinct check types minimizes miss rate at fixed cost, since none decomposes the budget into more than one channel to begin with. Neither of the two reproducible baseline arms in this paper's Experiments section (\S\ref{sec:experiments}) is compared against on its own published safety numbers --- the suspicion-score-triggered auditing \emph{protocol} established by \citet{AIControl2024} and \citet{ControlTax2024} is reproduced at matched call budget as one external comparison arm, not either paper's own reported result.

\textbf{Attacker-side adaptation against single-channel monitors.} A fast-moving line stresses those same protocols from the attacker's side, and it establishes the premise of this paper's RQ4 as demonstrated prior work rather than a hypothesis: adaptive attackers that know the protocol and the monitor consistently evade diverse LM monitors across AI-control benchmarks --- with resampling-based defenses actively backfiring \citep{AdaptiveAttacksMonitors2025} --- and strategic attack \emph{selection} (choosing when to attack or abort) cuts measured safety by 20 or more percentage points at roughly 1\% audit budgets without any new attack capability \citep{AttackSelectionConcentrated2026,AttackSelectionAgentic2026}. We concede that fragility of one predictable channel outright. What that line never varies is the defender's budget composition: every attack adapts against a single monitor signal, and no matched-cost comparison exists between that single channel and a portfolio of structurally distinct executable checks under the same aware adversary --- the specific gap RQ4 pre-registers, with the attack-selection results doubling as evidence that the allocation-aware condition this paper builds is a realistic, not paranoid, threat model.

\textbf{Formal inspection and security-game theory.} Classical inspection-game theory \citep{AvenhausInspectionGames2002} formalizes budget-constrained inspector-inspectee interactions and shows that optimal inspection under scarcity is typically a mixed strategy, the standard theoretical ancestor of this paper's framing; a more recent generalization of repeated principal-agent security games \citep{CalibratedStackelberg2024} characterizes the principal's achievable utility when the agent best-responds to calibrated forecasts of the principal's strategy rather than observing it directly, and AI-control deployment protocols themselves have been formalized as multi-objective audit games over a trusted monitor's signal \citep{GamesForAIControl2024} --- still with one homogeneous monitoring technology as the only inspection instrument. All of these give rigorous characterizations of achievable value under strategic pressure, but at the level of abstract payoff matrices or a single homogeneous inspection technology --- none instantiates ``inspection actions'' as concrete, structurally distinct verification operations on LLM-agent code, and none produces an empirical miss-rate-versus-cost frontier. A third, applied line --- deployed Stackelberg security games, which randomize the allocation of limited inspection or patrol resources against strategic attackers \citep{Tambe2011SecurityGameTheory,SinhaEtAl2018SSG} --- makes the directional content of RQ3 the expected result: that randomized or mixed allocations degrade less against an informed adversary than deterministic ones is close to a theorem in that literature. We concede the direction and position RQ3's contribution as the measured magnitude of that predicted effect in a regime the security-games line does not cover: executable check flavors with heterogeneous, computable detection profiles, a per-call cost meter over code-verification operations, and an empirical code-writing attacker.

\textbf{Reward-hacking and proxy-optimization mitigation.} One line of work \citep{InfTimeRewardHacking2024} formalizes the failure mode this paper's motivation depends on directly: reward hacking under inference-time optimization as a winner's-curse phenomenon, with a tuning procedure that hedges optimization strength against a fixed, imperfect proxy to avoid a peak-and-collapse in true reward. A complementary line \citep{CorrelatedProxies2024} gives a definition-level account of when a proxy is valid at all, showing that occupancy-measure regularization can provably prevent the optimization-induced breakdown of proxy validity. Both mitigate the same underlying strategic-exploitation dynamic this paper measures, but entirely from the optimizer's side, holding the verifier itself fixed; neither treats the complementary defender-side question of how a fixed budget of additional, independent checks should itself be allocated once the agent's optimization pressure is taken as given.

\textbf{Adversarial test-suite generation.} A training-time approach \citep{UTRL2024} jointly trains a unit-test generator and a code generator, rewarding test suites for discriminating against sampled imperfect programs drawn from the coder's own natural error distribution, producing an adaptive curriculum without gold test annotations. This establishes that test suites can be optimized against a moving target, but the target is a sampled distribution of natural mistakes rather than a strategic adversary that models the verifier and concentrates shortcuts where checks are absent, and the setting is training-time test generation rather than an evaluation-time budget split across check types --- the gap this paper's blind-versus-aware stress test is specifically built to close.

\textbf{Combinatorial test design.} A long-standing software-testing literature \citep{KuhnKackerLei2013CombinatorialTesting} designs covering arrays and pairwise/$t$-way interaction test suites to guarantee combinatorial input-space coverage against a fixed parameter model --- the construction technique this paper's hidden-combinatorial-test flavor is built directly on top of (a non-corpus canonical entry, cited here for the technique's provenance rather than as a retrieved paper in the search corpus). That literature optimizes coverage completeness against unstructured, non-adversarial faults under a static input-parameter specification; it has no notion of a strategic, pass-rate-optimizing agent that can concentrate implanted shortcuts specifically where a check budget does not reach, no cost meter for comparing one check flavor's expense against another's, and no allocation-aware attacker. It supplies the construction technique this paper's $H$ flavor borrows, without addressing --- or even being able to pose --- the composition question this paper asks.

\textbf{Metamorphic testing.} The invariance-probe flavor $P$ is built directly on the metamorphic-testing literature, which tests programs without output oracles by checking metamorphic relations --- semantics-preserving input transformations whose output consistency is machine-verifiable. Exactly as with combinatorial test design above, that literature supplies the construction technique only: it targets unstructured, non-adversarial faults under a fixed specification, with no strategic pass-rate-optimizing agent that can concentrate shortcuts where checks are absent, no cross-flavor cost meter, and no budget-composition question. It answers ``how do I check a program with no oracle'' \citep{SeguraEtAl2016MetamorphicSurvey,ChenEtAl2018MetamorphicSurvey}, not ``how should a fixed budget be split across structurally different checks against an adversary.''

\textbf{Backdoor construction technique provenance.} The three attack families used to implant machine-checkable shortcuts (\S\ref{sec:methodology}) build on established backdoor-construction techniques from adjacent domains: iterative trigger-injection strategies for textual backdoors \citep{acl_yan-etal-2023-bite} and input-conditional backdoor insertion into code-search models \citep{acl_sun-etal-2023-backdooring} both demonstrate that a machine-checkable, deliberately-implanted trigger can substitute for naturally occurring vulnerabilities without contaminating the surrounding classification/ranking task --- the same substitution principle this paper's construction oracle (\S\ref{sec:methodology}) relies on for coding tasks specifically.

\section*{Threats to Validity}

\textbf{Construct.} The central construct --- a genuine-versus-shortcut label $y_e$ --- is produced by construction (implanted witnesses) rather than observed. Its fidelity to shortcuts a real agent would discover unprompted, rather than ones deliberately implanted for this study, is an assumption, not a proof; the Problem Definition states this substitution explicitly rather than treating it as transparent, and a small validation audit against naturally occurring test-special-casing episodes (\S\ref{sec:experiments}, ``Further evaluation'') is the partial, disclosed mitigation.

\textbf{Internal.} The blind-versus-aware comparison (RQ3, RQ4) is only informative if the allocation-aware attack condition genuinely captures what a persistent, learning adversary could infer in deployment, rather than an artificially strong or artificially weak proxy for that adversary; the attack policy's information set (past check outcomes plus an allocation descriptor) is a modeling choice, stated explicitly, not an empirical measurement of real deployment leakage. A second internal threat is elicitation strength: if prompted attacker models refuse implantation or produce only weak shortcuts, every frontier flattens and flavor differences become unmeasurable for reasons unrelated to composition --- mitigated by the scripted-implantation floor and the open-weights attacker fallback (\S\ref{sec:experiments}), and detected by the walking-skeleton separation gate and RQ3's manipulation check rather than assumed away.

\textbf{External.} Every research question is scoped to implanted shortcuts in Python coding tasks (APPS, BigCodeBench, EvalPlus-extended HumanEval+/MBPP+) with three specific attack families; any allocation recommendation is conditional on this attack-family mix and this language/domain, not a claim about shortcut-taking in other agent domains (web agents, long-horizon SWE tasks) or under naturally emergent shortcuts, both disclosed as out of scope in the Conclusion.

\textbf{Conclusion validity.} All five RQs pre-register a Holm-corrected significance test and a task-level cluster bootstrap as the primary interval estimate, specifically to avoid inflating false-positive composition claims across the many pairwise flavor/policy/budget contrasts; the thin-cell rule (\S\ref{sec:experiments}) additionally guards against treating an underpowered cell's null result as evidence of no effect.

\section{Conclusion}

This work reframes a governance assumption behind budgeted verification --- that a fixed inspection budget's \emph{size} is the variable that matters --- as an empirical question about its \emph{composition}. The central methodological wager is that composition effects are only visible once structurally different checks are made commensurable under one cost meter, swept across the full allocation simplex rather than compared as isolated pure strategies, and tested against an attacker that can adapt to the verifier's own policy rather than only a naive one; a weaker design risks mistaking a naive-attacker artifact for genuine robustness, or mistaking held-out novelty for targeting structure. Whichever composition proves best is decision-relevant either way: if targeted flavors and mixtures hold up under an adaptive attacker, verifier designers gain a concrete recommendation for spending a fixed budget; if the advantage collapses under adaptive pressure or the novelty-ablation control, that is itself the finding --- the field's intuitions about ``smarter'' checks may be substituting novelty for real targeting power. Two limitations bound this claim from the outset: the study is scoped to implanted shortcuts in Python coding tasks, leaving naturally emergent shortcut-taking and other agent domains for follow-up; and any recommended allocation is conditional on the attack-family mix tested here, a scope this design reports explicitly rather than silently generalizing beyond.

\vspace{1em}
\hrule
\vspace{1em}

\section*{Pre-handoff checklist}
\begin{itemize}[leftmargin=1.4em]
\item \textbf{Organization/navigation} --- narrative flows problem $\to$ method $\to$ experiments $\to$ related work without forward references to undefined terms.
\item \textbf{Terminology discipline} --- ``check flavor,'' ``allocation policy,'' ``miss rate,'' ``allocation-blind/-aware'' each defined at first use in Problem Definition, used consistently thereafter; no internal stage IDs (\texttt{sp\_*}, \texttt{m\_*}, \texttt{rq\_*}, C1--C4, M1--M5) leaked into prose --- checked directly against the discussion pipeline log.
\item \textbf{Notation with worked example} --- the formal setup section defines $(\tau, r_\tau, s, e, y_e, w_e, F, \pi, B, MR)$ and the Problem Definition/Methodology/Experiments sections consistently reuse this notation; no separate worked numeric example is included since no numbers exist yet (would require fabricating results).
\item \textbf{RQ framing matches outcomes} --- all 5 RQs are graded/quantitative (rates, AUC differences, degradation gaps), each answered in its own subsection with an explicit falsifiable expected outcome, not a binary yes/no; RQ4 was specifically reworded during the discussion loop's Stage 6 refiner pass to avoid a ``beats baseline'' framing inconsistent with this paper's \texttt{study\_type: empirical} classification.
\item \textbf{No typesetting artifacts} --- only explicit placeholder markers in Results, per the drafting convention this is a pre-registered spec, not a report.
\item \textbf{Artifact plan} --- the verification-flavor library, attack-policy suite, and per-episode logs are the natural release artifacts; not yet built, so not claimed as available.
\end{itemize}

\section*{Hand-off}
This is a draft, not a finished paper --- experiments have not run. Natural next steps: prose polish and venue-specific LaTeX formatting once a target venue is chosen, or revising an earlier stage if piloting surfaces a problem (e.g., if the walking-skeleton gate shows the three flavors don't separate at all on scripted attacks).

\bibliographystyle{plainnat}
\bibliography{refs}

\end{document}
